\section{Propagation}
\label{sec:propagation}

This section transcribes \spec{9.2}, the three-phase procedure that assigns
each package a propagated channel, resolves every package's channel, and then
assigns propagated bumps. The phase order is forced: phase 3 reads what phase 2
produces, and \spec{9.2} forbids fusing them.

\subsection{The procedure, as written}

The pseudo-code is reproduced from \spec{9.2} unchanged, so that a reader can
diff it against \code{SPEC.md}. The definitions that follow are its
clauses, one by one.

\begin{speccode}[fontsize=\footnotesize]
propagate(units, graph, held, W):
    edges = graph.restrictedTo(config.propagation.kinds)          # §8.4

    # ---------- shared, computed once ----------
    # sourcePackages is a pure function of the unit (§13.4a) and both passes read it.
    # Computing it inside each pass evaluates it twice for every unit carrying both axes.
    src = { u: sourcePackages(u) for u in units }                 # §13.4a

    # ---------- phase 1: channel axis ----------
    chan = {}                                   # package -> (channel, commit, unitIndex)
    for u in units:
        if u.channelDepth == 0:                     continue      # §8.3a: default
        if u.propagateChannel == none:              continue
        sources = src[u]
        if sources is empty:                        continue
        c      = commitOf(u)                                      # loop-invariant
        pscope = resolve(u.propagateChannelScope)                 # §8.5a: resolved once,
        origin = UNCOMPUTED                                     # memoized on first admitted target
        for (d, level) in reach(edges, sources, u.channelDepth):  #   never per target
            if d not in pscope:                           continue
            if c not in Wfresh(d):                             continue # admission, §13.4a
            if cancelledFor(c, d):                        continue # §13.5a
            if u.propagateChannel == 'inherit' and origin == UNCOMPUTED:
                origin = originChannel(u)
            v = propagatedChannelFor(u, d, origin)                 # §9.3; origin computed once
            if v is NONE:                                 continue
            chan[d] = newerOf(chan[d], (v, u))                     # W160 on conflict

    # ---------- phase 2: channel resolution ----------
    channel = resolveChannels(chan, units)                         # §13.8

    # ---------- phase 3: bump axis ----------
    prop = {}                                   # package -> bump
    prov = {}                                   # package -> set of origin packages
    for u in units where bumpOf(u) != none:
        b = (u.propagate == 'inherit') ? bumpOf(u) : u.propagate
        if b == none or u.depth == 0:                continue
        sources = src[u]
        if sources is empty:                        continue
        c       = commitOf(u)
        pscope  = resolve(u.propagateScope)                        # §8.5: once per unit
        # §9.3a resolvability depends only on the unit's sources, so it is decided here,
        # once, and reduces to membership in the distinct source-channel set per target.
        srcChan   = { channel[P] : P in sources }                  # §9.3a
        anyStable = 'stable' in srcChan
        for (d, level) in reach(edges, sources, u.depth):
            if d not in pscope:                           continue
            owed = { P in sources : not delivered(P, c, d) }    # admission, §13.4a
            if owed is empty:                             continue
            if cancelledFor(c, d):                        continue
            if not (anyStable or channel[d] in srcChan):  continue # §9.3a, W208
            prop[d] = max(prop[d], b)
            prov[d] = prov[d] | owed
    return prop, prov, chan
\end{speccode}

The two axes admit differently, and the difference is the subject of this
report's revision. Phase 1 tests \code{c not in Wfresh(d)}: a dependent is
reached on the channel axis if and only if it has not itself released past the
commit carrying the unit. Phase 3 tests \code{owed is empty}: a dependent is
reached on the bump axis if and only if some source of the unit has not yet
\emph{delivered} that commit to it (\cref{def:delivery}). The two coincide
except for a dependent that released past the commit before its source did
(\cref{lem:deliv-window}), and \code{prov[d]} records exactly the owed
sources.

\subsection{The traversal}

\begin{definition}[reach, \spec{9.2}]
\label{def:reach}
For $S \subseteq \Pkgs$ and $\ell \in \mathbb{N}\cup\{\infty\}$,
\[
  \reachset{S}{\ell} \;=\;
  \{\, d \in \Pkgs \setminus S \;:\; \dist{S}{d} \leq \ell \,\},
\]
and each $d \in \reachset{S}{\ell}$ carries the level $\dist{S}{d}$.
\end{definition}

\begin{speccode}
reach(edges, sources, depth):
    frontier = sources
    seen     = set(sources)                     # a unit never reaches its own packages
    out      = {}                               # package -> level
    for level in 1 .. depth:                    # 'all' => until frontier is empty
        next = {}
        for p in frontier:
            for d in edges.dependentsOf(p):
                if d in seen: continue
                seen.add(d); next.add(d)
        if next is empty: break
        for d in next: out[d] = level
        frontier = next
    return out
\end{speccode}

\begin{lemma}[the traversal computes \cref{def:reach}]
\label{lem:reach}
On termination, \code{out} has domain $\reachset{S}{\ell}$ and
$\code{out}[d] = \dist{S}{d}$.
\end{lemma}

\begin{proof}
Standard breadth-first search from the set $S$, seeded at level $0$. By
induction on \code{level}: \code{frontier} at the start of iteration $n$ is
exactly $\{d : \dist{S}{d} = n-1\} \setminus S$ for $n>1$ and $S$ for $n=1$,
because \code{seen} contains precisely the packages at distance $\leq n-1$
(together with $S$ at level zero), and a package first enqueued at iteration $n$ has a chain of
length $n$ from $S$ and none shorter. Membership of \code{seen} is tested
before insertion and never cleared, so each package is assigned at most one
level, its least. The loop runs for $\ell$ iterations or until the frontier is
empty. Seeding $\code{seen} = S$ removes $S$ from the output, which is
\cref{def:reach}'s $\Pkgs \setminus S$.
\end{proof}

Three properties recorded by \spec{9.2} are visible in this statement and are
used repeatedly below. Depth is shortest-path, so a package reachable at both
depth $1$ and depth $3$ is included by \code{+1}. Depth is measured from the
originating source set as a whole and never recomputed from an intermediate
package, so a unit written \code{+1} reaches exactly the direct consumers of its
own packages in every run, whatever released in between. And a package reached
by several paths is visited once, so a diamond does not double-count.

\begin{example}[depth, on \cref{fig:workspace}]
\label{ex:reach}
Take $S = \{\code{core}\}$, the source set of a unit scoped \code{(core)}:

\begin{center}\small
\begin{tabular}{@{}L{2.6cm}L{5.4cm}L{4.6cm}@{}}
\toprule
Written & $\reachset{S}{\ell}$ & Note \\
\midrule
\code{feat(core)} & $\emptyset$ & $\ell = 0$, the default \\
\code{feat(core)\^{}} & $\{\code{cli},\code{ui},\code{api}\}$ & direct consumers \\
\code{feat(core)+2} & $\reachset{S}{1} \cup \{\code{docs-site},\code{@acme/theme}\}$ & two edges out \\
\code{feat(core)\^{}\^{}} & the same set & $\ell = \infty$; the graph is only two deep \\
\bottomrule
\end{tabular}
\end{center}
Now take the two-package source set $S' = \{\code{core},\code{ui}\}$, as a unit
scoped \code{(core,ui)} would have. Then $\reachset{S'}{1} =
\{\code{cli},\code{api},\code{docs-site},\code{@acme/theme}\}$: \code{ui} is
\emph{absent} although it is a dependent of \code{core} at distance $1$,
because \cref{def:reach} removes the whole of $S'$. The unit bumps \code{ui}
directly, and a package is never both a direct and a propagated target of the
same unit. \Cref{prop:bucket} is about exactly this clause.
\end{example}

\begin{remark}[self-exclusion is per unit]
\label{rem:selfex}
$\reachset{S}{\ell}$ excludes all of $S$, not merely the source a given path
started from. This is what stops a unit from propagating to a package it
already bumps directly. It is also the clause that makes the bucketed traversal
of \spec{13.11} unsound if transcribed naively; see \cref{prop:bucket}.
\end{remark}

\subsection{Phase 1: the channel axis}

\begin{definition}[proposed channel, \spec{9.3}]
\label{def:propchan}
For a unit $u$ and package $d$, let $cur = \chanof{\base{\st}{d}}$ (\stable{}
if $d$ has no baseline). Then $\mathrm{pcf}_{\st}(u,d) \in \Chans \cup
\{\NONE\}$ is:
\begin{itemize}
  \item if $\pchan{u}$ is a transition $\langle f \to t \rangle$: \NONE{} if
    $f = t$ (\diag{W207}); \NONE{} if $cur$ does not match $f$, where $f = *$
    matches any $cur \neq \stable$ and otherwise $f$ matches only $cur = f$;
    else $target = t$;
  \item otherwise $target$ is $\pchan{u}$, or the origin channel of $u$ where
    that is \code{inherit}; and if $target = \stable$ while $cur \neq \stable$,
    the result is \NONE{} (\diag{W200}: a propagated channel never graduates
    implicitly);
  \item in either case, if $target = cur$ the result is \NONE{} (\diag{W199});
    otherwise it is $target$.
\end{itemize}
The origin channel of $u$ is the channel $u$'s own \code{Channel} directive
assigns to its sources if it carries one, and otherwise
$\chanof{\base{\st}{p}}$ for its sources $p$, taking the value of the
byte-wise least name if they disagree (\diag{W160}).
\end{definition}

The final clause complements the target \Winf{} admission gate: it compares
$target$ against the package's \emph{baseline} channel rather than against
anything computed in the current run. This no-op check is the other premise of
\G{7}, and \cref{thm:g7} is its consequence.

\begin{definition}[phase 1 output, \spec{9.2}]
\label{def:phase1}
$\mathrm{prop}\chi_{\st}(d)$ is the value $\mathrm{pcf}_{\st}(u,d)$ of the
$\prec$-greatest unit $u$ (newest commit first, then last unit within a commit,
\spec{11.6}) satisfying: $\pcdepth{u} \neq 0$; $\pchan{u} \neq \code{none}$;
$\srcp{\st}{u} \neq \emptyset$; $d \in \reachset{\srcp{\st}{u}}{\pcdepth{u}}$;
$d \in \pcscope{u}$; $\cof{u} \in \Winf{\st}{d}$; $\lnot\canc{\cof{u}}{d}$; and
$\mathrm{pcf}_{\st}(u,d) \neq \NONE$. It is undefined if no unit qualifies.
\end{definition}

The channel axis keeps $\cof{u} \in \Winf{\st}{d}$ as its admission because a
channel is carried by the units and needs no release of the source
(\spec{13.4a}, \G{7}); it has no notion of delivery.

\subsection{Phase 2: channel resolution}

\begin{definition}[channel resolution, \spec{13.8}]
\label{def:channels}
For every $p \in \Pkgs$, whether or not it will be released, let $base =
\chanof{\base{\st}{p}}$ and let $D_{\st}(p)$ be the value proposed by the
$\prec$-greatest unit $u$ with $p \in \scopef{u}$, $\cof{u} \in \Winf{\st}{p}$,
$\lnot\canc{\cof{u}}{p}$, that carries a \code{Channel} directive proposing
something for $p$ (a transition matching $base$, or a value differing from
$base$; a directive proposing nothing is not a competitor and does not count
towards \diag{W186}). Then
\[
  \chanres{\st}{p} \;=\;
  \begin{cases}
    D_{\st}(p) & \text{if defined,}\\
    \mathrm{prop}\chi_{\st}(p) & \text{else if defined,}\\
    base & \text{otherwise.}
  \end{cases}
\]
\end{definition}

A direct directive beats every propagated one regardless of age; among direct
directives, and separately among propagated ones, the newest commit wins and
then the last unit within it. A package named by nothing keeps its baseline
channel, which is why an established train needs no directives to stay
together.

\subsection{Phase 3: the bump axis}

\begin{definition}[resolvability, \spec{9.3a}]
\label{def:resolvable}
$\srcch{\st}{u} = \{\chanres{\st}{p} : p \in \srcp{\st}{u}\}$ and
\[
  \resolv{\st}{u}{d} \;\equiv\;
  \stable \in \srcch{\st}{u} \;\lor\; \chanres{\st}{d} \in \srcch{\st}{u}.
\]
\end{definition}

A bump is a claim that the dependent has something new to pick up, and across a
channel boundary the claim is false: a stable consumer resolves a
beta-releasing dependency by its stable range exactly as before, so a
republished consumer would contain byte-for-byte what its predecessor
contained. Suppression is reported per suppressed pair as \diag{W208}. Note
that \code{srcChan} in the pseudo-code is computed over the whole source set,
not over the owed sources; \cref{sec:duality} returns to this.

\begin{example}[a caret that reaches and does not bump: \SPEC{} vector 95]
\label{ex:resolvable}
The unit \code{feat(core)\^{}\%beta: x} on \cref{fig:workspace}. The \code{\%}
puts \code{core} on \code{beta}; the \code{\^{}} reaches
$\{\code{cli},\code{ui},\code{api}\}$. Phase 2 leaves the three consumers where
their baselines put them, on \stable, so
\[
  \srcch{\st}{u} = \{\code{beta}\},\qquad
  \stable \notin \srcch{\st}{u},\qquad
  \chanres{\st}{\code{cli}} = \stable \notin \srcch{\st}{u},
\]
and \cref{a:resolvable} fails for every one of them: three \diag{W208}
diagnostics, and $\plan{\st} = \{\code{core}\}$ alone. The caret was honoured
and reached its targets; what it would have delivered is not resolvable by
them. Writing \code{feat(core)\^{}\%beta++1} instead moves the same three onto
\code{beta} in phase 1, so $\chanres{\st}{\code{cli}} = \code{beta} \in
\srcch{\st}{u}$ and all four packages release, which is the form that keeps a
train installable as a set.
\end{example}

\begin{definition}[admission, \spec{9.2}, \spec{13.4a}]
\label{def:admission}
For $u \in \Units$ and $d \in \Pkgs$, $\admit{\st}{u}{d}$ holds iff all of:
\begin{enumerate}[label=(A\arabic*),leftmargin=3.2em]
  \item $\bumpof{u} \neq \bnone$, $\pbump{u} \neq \bnone$, $\pdepth{u} \neq 0$;
    \label{a:live}
  \item $\srcp{\st}{u} \neq \emptyset$ and
    $d \in \reachset{\srcp{\st}{u}}{\pdepth{u}}$; \label{a:reach}
  \item $d \in \pscope{u}$; \label{a:scope}
  \item $\owed{\st}{u}{d} \neq \emptyset$; \label{a:window}
  \item $\lnot\canc{\cof{u}}{d}$; \label{a:cancel}
  \item $\resolv{\st}{u}{d}$. \label{a:resolvable}
\end{enumerate}
\end{definition}

\begin{definition}[propagated and direct bumps, \spec{13.6}, \spec{13.7}]
\label{def:propagated}
\label{def:direct}
\begin{align*}
  \propb{\st}{d} &= \bigsqcup \{\, \pbump{u} \;:\; u \in \Units,\ \admit{\st}{u}{d} \,\}, \\
  \dirb{\st}{p} &= \bigsqcup \{\, \bumpof{u} \;:\; p \in \srcp{\st}{u} \cap \scopef{u},\
     \cof{u} \in \Winf{\st}{p} \,\},
\end{align*}
both empty joins being \bnone. Then $\effb{\st}{p} = \dirb{\st}{p} \sqcup
\propb{\st}{p}$.
\end{definition}

\Cref{a:window} is the load-bearing clause of the entire design, and it is
worth being explicit about what it is and is not. It tests the \emph{target's}
position against its \emph{sources'} releases: $d$ is admitted while some
source has not delivered $\cof{u}$ to it. It does not test the source's window,
and it does not test the target's window either. \spec{13.7a} shows what each
of the other readings costs. With retention governed by the source's window, a
dependent whose publish failed while its provider's succeeded is silently
orphaned, permanently one patch behind a dependency it declares, with no
diagnostic, because from the engine's point of view nothing is pending. With
admission governed by the target's window, the rule of CCME 2.0.0, the same
orphan is reached from the other side: a dependent that proceeded past the
commit on a cause of its own while its provider failed is owed nothing once
the provider publishes, and keeps the range it reconciled against the old
version. By \cref{lem:deliv-window} the delivery test agrees with the window
test on every target that has not released past $\cof{u}$; it differs only on
a target that has, where it asks the finer question. Every guarantee in
\cref{sec:guarantees} rests on \cref{a:window} being this test, and
\cref{thm:g2} is where that is used.

\begin{example}[the orphan, from both sides: \SPEC{} vectors 66 and 80d]
\label{ex:orphan}
Both shapes use the history of \cref{ex:window}, $c_0 \to c_1 \to c_2 =
\Head$, with the tags of $\st_0$ at $c_0$ and the unit $u$ =
\code{feat(core)\^{}: streaming reader} at $C = c_2$. So $\scopef{u} =
\{\code{core}\}$, $\bumpof{u} = \bminor$, $\pbump{u} = \bpatch$ (the default of
\cref{def:unit}) and $\pdepth{u} = 1$.

\emph{Provider fails to deliver; the consumer waits} (vector 66). Run 1
evaluates \cref{def:admission} with $\st_0$. All of \code{cli}, \code{ui},
\code{api} lie in $\reachset{\{\code{core}\}}{1}$, none is cancelled, all are
on \stable{} as is \code{core}, and $C \in \Winf{\st_0}{d}$ for each, so by
\cref{lem:deliv-window} $\owed{\st_0}{u}{d} = \{\code{core}\}$, all three are
admitted and
\[
  \plan{\st_0} = \{\code{core},\code{cli},\code{ui},\code{api}\},\qquad
  \vers{\st_0}{\code{core}} = 1.5.0,\quad
  \vers{\st_0}{\code{cli}} = 2.0.1 .
\]
Suppose \code{core}, \code{ui} and \code{api} publish and \code{cli} fails, so
\[
  \st_1 = \st_0 \cup \{\code{core@1.5.0},\ \code{ui@0.9.2},\ \code{api@1.2.1}\}
  \quad\text{at } C.
\]
Run 2 is where the readings of \cref{a:window} separate. Under delivery, the
normative reading, $C \not\anc \basec{\st_1}{\code{cli}} = c_0$, so
$\lnot\deliv{\st_1}{\code{core}}{C}{\code{cli}}$ and
$\owed{\st_1}{u}{\code{cli}} = \{\code{core}\}$: \code{cli} is admitted. Under
the target's window, the reading of CCME 2.0.0, $C \in \Winf{\st_1}{\code{cli}}
= \{c_1, C\}$: \code{cli} is admitted. Under the source's window, the trap,
$C \notin \Win{\st_1}{\code{core}} = \emptyset$: $u$ reaches nobody.
The first two agree, as \cref{lem:deliv-window} says they must for a
\code{cli} that has not released past $C$. For \code{ui}, which has,
$\deliv{\st_1}{\code{core}}{C}{\code{ui}}$ holds with $t = C \anc C =
\basec{\st_1}{\code{ui}}$: nothing is owed, and \code{ui} is out. Under the
normative rule $\plan{\st_1} = \{\code{cli}\}$ at the version run 1 planned,
marked \diag{W193}, and $\plan{\st_2} = \emptyset$. Under the trap
$\plan{\st_1} = \emptyset$ and \code{cli} is never released, on this run or any
later one, with no diagnostic anywhere.

\emph{Consumer gets ahead of its provider} (vector 80d). Let $C$ carry a
second unit $u'$ = \code{feat(cli): y} beside $u$. Then $\dirb{\st_0}{\code{cli}}
= \bminor$, and run 1 plans $\plan{\st_0} =
\{\code{core},\code{cli},\code{ui},\code{api}\}$ with $\vers{\st_0}{\code{cli}}
= 2.1.0$. Suppose \code{core} fails. \code{ui} and \code{api} have no cause of
their own relative to $X = \{\code{core}\}$ and are blocked (\diag{W194});
\code{cli} has one, its direct minor, and \emph{proceeds}
(\cref{def:owncause}): it is published at $2.1.0$ with its manifest reconciled
to \code{core}'s published baseline $1.4.2$, and tagged at $C$. So
\[
  \st_1 = \st_0 \cup \{\code{cli@2.1.0}\ \text{at}\ C\},\qquad
  \Winf{\st_1}{\code{cli}} = \emptyset .
\]
Run 2. \code{core} is planned by its direct minor. For \code{cli}, the window
reading says $C \notin \Winf{\st_1}{\code{cli}}$ and admits nothing: \code{cli}
keeps the range it reconciled against $1.4.2$, is never planned again, and is
one minor behind a dependency it declares, silently. The delivery reading
asks whether \code{core} has delivered $C$ to \code{cli}: $\relc{\st_1}{\code{core}}
= \{c_0\}$ and $C \not\anc c_0$, so $\owed{\st_1}{u}{\code{cli}} =
\{\code{core}\}$, \code{cli} is admitted with $\propb{\st_1}{\code{cli}} =
\bpatch$, and
\[
  \plan{\st_1} = \{\code{core},\code{cli},\code{ui},\code{api}\},\qquad
  \vers{\st_1}{\code{cli}} = 2.1.1,
\]
ordered after \code{core} and blocked if \code{core} fails again. $\st_1$ is
well-formed (\cref{def:wf}): \code{core} would be released at $C =
\basec{\st_1}{\code{cli}}$ and still owes \code{cli}, and \code{cli} is in the
plan after it. Once both are tagged, $\deliv{\st_2}{\code{core}}{C}{\code{cli}}$
holds with $t = C$, nothing is owed, and $\plan{\st_2} = \emptyset$. The plan
of run 2 equals the plan of run 1 although run 1 published a package; that is
the observation \cref{thm:g6} has to accommodate.
\end{example}

\begin{example}[two tags on one commit: \diag{E201}, vector 80d continued]
\label{ex:e201}
Continue the second shape from $\st_1$. Suppose run 2 publishes \code{core},
\code{api} and \code{ui} and \code{cli} fails. Then $\st_2 = \st_1 \cup
\{\code{core@1.5.0}, \code{api@1.2.1}, \code{ui@0.9.2}\}$ at $C$, and now
$\deliv{\st_2}{\code{core}}{C}{\code{cli}}$ holds with $t = C \anc C =
\basec{\st_2}{\code{cli}}$: the ledger reads \code{cli} as served, although
\code{cli@2.1.1} was never published and \code{cli@2.1.0} still names
$1.4.2$. $\owed{\st_2}{u}{\code{cli}} = \emptyset$, $\plan{\st_2} = \emptyset$,
and no later plan at any commit will compute the debt, because two tags on one
commit have no order. This is \cref{prop:loss}: run 2 reports \diag{E201}
for the pair, and the remedy is an explicit \code{Release-As} on \code{cli}
at a later commit. The other route to the same state is prevented rather
than reported: a run 2 in which \code{cli} sat out, held or deselected, would
release \code{core} at \code{cli}'s baseline commit without \code{cli}
following, and \spec{19.3} refuses it before any publication. Had \code{core}
instead published at a later commit $c_3$ in a run \code{cli} sat out, $C$
would stay in $\Wowed{}{\code{core}}{\code{cli}} = \Rch{c_3}\setminus\Rch{c_0}$
(\cref{def:owedwindow}) and the next run would plan \code{cli}'s catch-up.
\end{example}

\begin{remark}[catch-up is not a mode]
There is no separate catch-up pass in this model, and none in the
specification. Catch-up is what the algorithm does when the ordinary rule is
evaluated against the right position (\spec{13.7a}). A ``catch-up release'' is
a label for a release whose entire cause is a propagation from a package not
itself in the run's plan (\diag{W193}); the engine does nothing different to
produce it, and vector 80d's second run is not one, since \code{core} is in
that plan. Catch-up does not widen depth, resurrect a cancelled contribution,
release a held package, change a planned version, or propagate onward from the
catching-up package; each of which is a property of \cref{def:admission}
rather than a rule about catch-up.
\end{remark}
